\chapter{Investigación de Operaciones: repartir el dinero}
\label{cap:s-dinero}
\textit{Materia: Investigación de Operaciones. Pregunta: ¿cuánto dinero va a cada lugar, en qué mes y por qué red social,
para traer a la mayor cantidad de gente sin llenar ningún lugar?}

\section{El problema, como una receta}
Todo problema de optimización tiene tres ingredientes:
\begin{enumerate}
\item \textbf{Lo que se decide} (las \emph{variables}): cuántos pesos poner en cada mes, en cada lugar y en cada canal.
  Son 9 meses (octubre de 2026 a junio de 2027) $\times$ 3 lugares (Chetumal, Bahía y Ruta) $\times$ 2 canales (Google y
  Facebook) = \textbf{54 cantidades}.
\item \textbf{Lo que se quiere lograr} (el \emph{objetivo}): el mayor número de visitantes esperados hacia el sur.
\item \textbf{Las reglas que no se pueden romper} (las \emph{restricciones}).
\end{enumerate}
La computadora busca, entre todas las formas posibles de repartir el dinero, la que logra el objetivo sin romper ninguna
regla.

\begin{porque}
\textbf{¿Por qué visitantes y no dinero que deja el turismo (derrama)?} La derrama que publica el gobierno no dice en qué
unidad está y termina en marzo de 2024; habría que inventar cuánto gasta cada turista. \textbf{¿Por qué no ``visitantes
menos un castigo por presión''?} Habría que inventar cuánto pesa el castigo. Brandon eligió \textbf{visitantes al sur}, y la
presión se controla con reglas duras.
\end{porque}

\section{Cuánto rinde un peso en cada red}
\begin{intuicion}
Un anuncio funciona como un embudo: mucha gente lo ve, algunos le dan clic y, de esos, algunos planean el viaje. Se paga
por cada clic. Lo que importa es cuántos viajeros trae cada peso.
\end{intuicion}

\begin{ejemplo}[ · Google contra Facebook]
Con los costos publicados para anuncios de viajes (promedios de Estados Unidos, usados como supuesto) y el dólar a 17.06
pesos:
\begin{itemize}
\item \textbf{Google}: un clic cuesta 2.12 dólares, es decir, $2.12\times17.06=36.17$ pesos. De cada 100 clics, 5.75 planean
  el viaje. Con \$1,000 se compran $1{,}000\div36.17=27.6$ clics y llegan $27.6\times0.0575=\mathbf{1.59}$ visitantes.
\item \textbf{Facebook}: un clic cuesta 0.51 dólares, 8.70 pesos. Con \$1,000 se compran 114.9 clics y llegan
  $\mathbf{6.61}$ visitantes.
\end{itemize}
Facebook rinde cuatro veces más por peso porque su clic es cuatro veces más barato.
\end{ejemplo}

\textbf{Un hueco importante}: nadie publica cuántas personas planean un viaje después de dar clic en un anuncio de turismo
en Facebook. El caso base usa la cifra de Google (5.75\,\%) y se prueba qué pasa con 3\,\% y con 6.38\,\%.

\section{Las cuatro reglas de cuidado}
\begin{center}
\small
\begin{tabularx}{\linewidth}{>{\raggedright\arraybackslash}p{0.3\linewidth}Y}
\encabezado \h{Regla} & \h{Por qué existe} \\
\textbf{Cero anuncio en temporada alta} & Anunciar un lugar cuando ya se llena sería crear la saturación que la campaña
quiere evitar. \\
\filagris\textbf{Tope de capacidad probada} & Los visitantes esperados más los de la campaña no pueden pasar del mes más
lleno que el lugar ha recibido en su historia. \\
\textbf{Mínimo de 15\,\% por lugar} & Para que el dinero no se vaya todo a un solo lugar y los tres reciban algo (equidad
entre comunidades). \\
\filagris\textbf{Tope de 70\,\% por canal} & Para no depender de una sola plataforma: si Facebook cambia sus precios o
sus reglas, la campaña no se cae. \\
\bottomrule
\end{tabularx}
\normalsize
\end{center}

\section{Lo ``estocástico'': decidir hoy sin saber cómo saldrá el año}
El dinero se decide hoy, pero el año puede salir malo, probable o bueno (los escenarios del Monte Carlo). Si sale bueno y
un lugar se llena, el anuncio de ese mes \textbf{se pausa} y esos visitantes se pierden. El modelo considera los tres
escenarios a la vez (con probabilidades de 30\,\%, 40\,\% y 30\,\%) y aprende, desde hoy, a no gastar en meses donde un año
bueno llenaría el lugar. A esto se le llama \textbf{programación estocástica de dos etapas}: la primera etapa es el plan;
la segunda, la reacción cuando se sabe qué pasó.

\section{Cómo lo resuelve la computadora}
\begin{intuicion}
Imagina todas las formas posibles de repartir el dinero como el interior de un diamante con muchas caras. Las reglas son
las paredes. La mejor solución siempre está en una de las esquinas. El método \textbf{simplex} camina de esquina en esquina
por las orillas, siempre hacia una mejor, hasta que ya no hay esquina vecina que mejore. Como pausar o no pausar es un sí o
no (no se puede pausar ``a medias''), se agrega otro método, \textbf{ramificación y acotamiento}, que prueba las dos
opciones y descarta los caminos que no pueden ganar.
\end{intuicion}
El programa PuLP escribe el modelo y el programa CBC lo resuelve en \textbf{0.3 segundos}.

\section{El resultado}
\begin{ejemplo}[ · el máximo de visitantes, a mano]
De los \$250,000 al año, a los 9 meses del plan les tocan $250{,}000\times9/12=\$187{,}500$.

Como Facebook rinde más, se lleva su tope de 70\,\%: \$131,250. Google, el resto: \$56,250.
\begin{itemize}
\item Facebook: $131{,}250\div8.70=15{,}084$ clics; $15{,}084\times0.0575=867$ visitantes.
\item Google: $56{,}250\div36.17=1{,}555$ clics; $1{,}555\times0.0575=89$ visitantes.
\end{itemize}
Total: \textbf{957 visitantes}, a unos \textbf{\$196} cada uno ($187{,}500\div957$).
\end{ejemplo}

\textbf{Reparto por lugar}: la Ruta, 41.5\,\% (\$77,790); la Bahía, 36.9\,\% (\$69,175); Chetumal, 21.6\,\% (\$40,535).
\textbf{Meses sin anuncio}: diciembre en los tres lugares, enero en la Ruta y la Bahía, marzo en la Ruta y abril en la
Bahía. Todos son de temporada alta: de 27 combinaciones de mes y lugar, 20 quedan permitidas.

\begin{error}
\textbf{El primer resultado mandaba todo a un solo mes.} Como no hay dato de que un lugar responda mejor que otro a los
anuncios, había muchísimas formas de repartir que daban los mismos 957 visitantes, y la computadora entregaba una
cualquiera: los \$131,250 de la Ruta, todos en junio. Matemáticamente era óptimo; en la práctica, absurdo. Brandon eligió
desempatar \textbf{en proporción al espacio libre de cada mes}: primero se calcula el máximo de visitantes, y luego, entre
las soluciones que lo logran, se elige la que reparte el dinero según cuánto espacio sobra.

Ejemplo: en octubre de 2026 se esperan 3,688 visitantes en la Ruta, que ha llegado a recibir 10,465. Le sobra el 64.8\,\%
de espacio. Sumando el espacio de las 20 combinaciones permitidas, a la Ruta en octubre le toca el 8.1\,\% del dinero de
cada canal: \$10,624 de Facebook y \$4,553 de Google.
\end{error}

\section{¿Cuánto cuesta cuidar?}
Para saber cuánto cuesta cada regla, se resuelve el modelo sin ella y se compara:
\begin{center}
\small
\begin{tabular}{lr}
\encabezado \h{Regla} & \h{Visitantes que cuesta} \\
Cero anuncio en temporada alta & \textbf{0} \\
Tope de capacidad probada & \textbf{0} \\
Mínimo de 15\,\% por lugar & \textbf{0} \\
Tope de 70\,\% por canal & 282 (29.5\,\%) \\
\bottomrule
\end{tabular}
\normalsize
\end{center}
\textbf{El hallazgo más bonito del proyecto}: las reglas de cuidado ambiental y de equidad \textbf{no le cuestan ni un
visitante} a la campaña. La campaña es chica frente al espacio que tiene el sur. La única regla que cuesta es el tope por
canal, y se mantiene por prudencia.

\textbf{Precios sombra}: la teoría de la optimización da, para cada regla, cuánto mejoraría el resultado si se aflojara un
poquito. Si al presupuesto se le agregan \$1,000, llegan 1.59 visitantes más (iría a Google, porque Facebook ya está en
su tope). Si a Facebook se le permiten \$1,000 más (quitándoselos a Google), llegan 5.02 más. Las reglas que no estorban
tienen precio sombra cero.

\section{¿Hasta dónde se puede apretar?}
Si se exige que ningún mes con anuncio pase de cierto porcentaje de su capacidad probada:
\begin{center}
\small
\begin{tabular}{lrrr}
\encabezado \h{Límite} & \h{100\,\% a 40\,\%} & \h{35\,\%} & \h{30\,\%} \\
Visitantes & 957 & 537 & 0 \\
\end{tabular}
\normalsize
\end{center}
La campaña puede prometer que \textbf{ningún mes con anuncio pasará del 40\,\% de lo que el lugar ya ha recibido}, sin
perder un solo visitante. A esto se le llama \emph{frontera de Pareto}: la línea donde ya no se puede mejorar una cosa sin
empeorar la otra.

\section{¿Y si los supuestos están mal?}
\begin{center}
\small
\begin{tabular}{lr}
\encabezado \h{Si cambia\dots} & \h{Visitantes} \\
Nada (caso base) & 957 \\
Facebook convierte solo 3\,\% & 542 \\
Facebook convierte 6.38\,\% & 1,052 \\
El clic de Facebook cuesta 0.42 dólares & 1,143 \\
Una tormenta quita 25\,\% o 50\,\% & 957 (los meses con riesgo ya están fuera) \\
\$500,000 al año & 1,914 \\
\$1,000,000 al año & 3,827 \\
\bottomrule
\end{tabular}
\normalsize
\end{center}
El número de visitantes cambia, pero \textbf{el reparto entre lugares y canales casi no cambia}. La recomendación es
estable aunque los supuestos se muevan.
